% NOT CURRENT.  The long form of the Prolegomena of version 12, as it stood on 25 November 2026
% after review round 9 and before it was shortened on the advice of rounds 8 and 9.
% The paper uses paper/v12_s0_prolegomena.tex.  This file holds the passages that were cut:
% the sphere-versus-cube comparison, the "machine with one dial", "The question of the minimum",
% "What the results give", and the closing paragraph on the two circles through the set of Peres.
% It is not compiled and its labels duplicate those of the current file.

\section{Prolegomena}\label{sec:prolegomena}

This section sets out the idea from which the paper grows and fixes the terms in which its
results are stated.  It contains no theorem.  A reader who wants the results first will find
them in Section~\ref{sec:intro} and can return here for the terms.

\subsection{A ray is a point on a cube}

Take first a ray with real coordinates: a straight line through a fixed point~$O$.  It has a
direction but no length and no sense.  Put $O$ at the centre of a cube.  A ray leaves the cube
at a point $P$ and at the opposite point~$P'$, and either point names the ray.  A set of rays,
which is hard to draw, becomes a set of points on the surface of a cube, each opposite pair
counted once.  (Rays with complex coordinates enter in Section~\ref{sec:grid}; they have no
such picture.)

A sphere would serve that purpose equally well.  The cube does one thing more.  Let its edges
run along three perpendicular axes and let it reach from $-x$ to $x$ on each, with $x>0$.  A
ray is a triple of numbers taken up to a common factor.  Choose the factor so that the
coordinate of largest size has size~$x$: the triple is then the point~$P$, and its coordinates
are read straight off the face on which $P$ lies.  Thus $(1,0,1)$ and $(x,0,x)$ are the same
ray, and the second is the point at which it meets the cube.  The sphere fixes the common
factor by the sum of the squares of the coordinates, the cube by the largest coordinate alone.
It is the cube's way that keeps simple coordinates simple.

\subsection{The rule}\label{sec:rule}

Three mutually orthogonal rays form a \emph{triad}.  The rule is:

\begin{quote}
In every triad exactly one ray is \emph{colored}.  Two orthogonal rays are never both colored.
\end{quote}

A ray that is not colored is \emph{uncolored}.  To \emph{color a set} of rays is to decide, for
every ray in it, whether it is colored or uncolored, in agreement with the rule.  For a finite
set the rule is applied to the orthogonal pairs and the triads found within the set, and to no
others.  The second sentence of the rule is then not idle.  Two orthogonal rays of the set
belong to a triad of space whose third ray may be absent from the set, and they are forbidden
to be both colored all the same.

The rule is a prediction of quantum mechanics.  For a particle of spin one, the squares of the
spin along three mutually orthogonal directions can be measured together.  Each square is
found to be 0 or~1 (in units of $\hbar^2$), and the three add up to~2.  So in every triad of
directions exactly one square is~0, and that direction is the one we color.  Suppose the
particle carried these values before the measurement, and suppose the value along a direction
did not depend on which two other directions were measured with it.  Then every ray in space
would be colored or uncolored once and for all, in agreement with the rule.

Kochen and Specker~\cite{KochenSpecker1967} showed that this cannot be done, and that one need
not consider every ray in space: there are finite sets of rays that cannot be colored.  Such a
set is \emph{uncolorable}, or a \emph{Kochen--Specker (KS) set}.  What it excludes should be
stated exactly: values that are fixed beforehand \emph{and} do not depend on what else is
measured.  A theory in which the value along a direction depends on the other two directions
chosen is not touched by it.

\subsection{The grid}\label{sec:grid}

Peres drew his set of 33 rays as points on a cube, and he drew the 31 rays of Conway and
Kochen the same way~\cite{Peres1991,Peres1993}.  In both drawings the points do not lie just
anywhere on the faces.  They lie where five lines in one direction cross five lines in the
other.  In the first drawing the lines sit at $0$, $\pm1$ and $\pm\sqrt2$.  In the second they
sit at $0$, $\pm1$ and~$\pm2$.

This is the observation on which the paper rests.  The two sets are built on \emph{the same
grid}, except for the position of one pair of lines.  Write that position as~$x$
(Fig.~\ref{fig:generic}).

\begin{figure}[t]
\centering
\includegraphics[width=0.6\textwidth]{figures/alphabet_cube_generic.pdf}
\caption{Coordinate triples as rays, for real $x>1$.  The cube reaches from $-x$ to $x$ on each
axis and the lines on its faces sit at $0,\pm1,\pm x$.  A bar is a minus sign.  Triples written
without~$x$, such as $(1,0,0)$ and $(1,\bar1,1)$, reach the surface after scaling by~$x$.  The
49 rays are named by the same coordinate triples for every such~$x$; which pairs are orthogonal
depends on~$x$.}
\label{fig:generic}
\end{figure}

The numbers at which the lines sit are the \emph{coordinate values}.  A \emph{coordinate
triple} is a triple of coordinate values, not all zero, such as $(x,1,\bar1)$, where a bar is a
minus sign.  Every coordinate triple names a ray, and proportional triples name the same ray.
The \emph{grid} is the set of rays so named.  Where a coordinate value is a parameter, as $x$
is here, a coordinate triple is an expression in the parameter, to be evaluated at each value
of it, and ``the same triple'' at two positions of the lines means the same expression.  For
real $x>1$ the coordinate values $0,\pm1,\pm x$ give 124 coordinate triples, and they name 49
rays.  Each of the 49 leaves the cube at a point where the lines cross, and every such point
is one of the 49: the drawing is the grid.  (For other real or complex $x$ the triples again
name 49 rays, except at a few values such as $x=1$, where fewer remain.)

The same terms are used for any finite list of numbers that contains~$0$.  The numbers may be
complex, although the cube can then no longer be drawn.  In counting coordinate values we
leave out $0$ and take a value and its negative as one, so the grid above has two coordinate
values, $1$ and~$x$.  Since a common factor does not change a ray, one of the nonzero coordinate
values can be taken to be~$1$.

Two orthogonal rays fix a third, the ray orthogonal to both, and it need not belong to the
grid.  With coordinate values $0,\pm1,\pm2$ the triples $(1,1,2)$ and $(2,0,\bar1)$ are
orthogonal, and the ray orthogonal to both is $(1,\bar5,2)$, which no coordinate triple names.
The \emph{completed grid} is the smallest set of rays that contains the grid and contains,
with every two orthogonal rays, the third ray of their triad.  A grid and its completed grid
are different sets, and every count in this paper says which is meant.

Earlier versions of this paper, and the programs that accompany it, use other names.  There
the list of coordinate values is an \emph{alphabet}, a coordinate value is a \emph{letter} or
\emph{symbol}, a coordinate triple is a \emph{word}, a grid is a \emph{pool} or \emph{raw
pool}, and a completed grid is a \emph{completed pool}.

\subsection{Orthogonality depends on where the lines sit}

The coordinate triples, as expressions, do not change with~$x$.  What changes is which rays
are orthogonal.
Two rays are orthogonal when the products of their coordinates add up to zero (with complex
coordinates one factor of each product is conjugated).  Take $x$ real and the triples
$(x,1,1)$ and $(x,\bar1,\bar1)$:
\[
(x,1,1)\cdot(x,\bar1,\bar1)=x^2-1-1=x^2-2 .
\]
These two rays are orthogonal when the lines sit at $\pm\sqrt2$ and at no other position.
Take instead $(1,\bar1,1)$ and $(x,1,\bar1)$:
\[
(1,\bar1,1)\cdot(x,1,\bar1)=x-1-1=x-2 .
\]
These are orthogonal when $x=2$ and not otherwise.

A warning is needed.  \emph{The picture shows where the rays are; it does not show which are
orthogonal.}  The cube with lines at $\sqrt2$ looks like the cube with lines a little to one
side, and no drawing can show that two rays are exactly orthogonal.  That has to be computed.
The picture's service is a different one.  It shows that there is a single quantity, the
position $x$ of a line, on which every orthogonality depends, and it invites us to slide that
line and watch what happens.

The computation is always of the kind just made.  The product of two coordinate triples is a
sum of at most three terms, and each term is a product of two coordinate values.  Either it
vanishes for every~$x$, as it does for the three axes, or it is an equation for~$x$ with small
whole numbers as coefficients, like the two above.  There are finitely many such equations.
Table~\ref{tab:zerosums} lists, up to the symmetries stated there, those with a solution other
than $0$ and~$\pm1$.  It is at the solutions of these equations that orthogonal pairs are
gained.  For real $x$ an equation that is not an identity has at most two solutions, so the
positions at which anything changes are finite in number.  More orthogonal pairs put more
constraints on a coloring.  Whether they leave none has to be computed in each case.

\subsection{The question}

We now have a machine with one dial.  Turn the dial and the coordinate triples stay the same
while the orthogonalities come and go (Fig.~\ref{fig:triptych}).  The question asks itself:

\begin{quote}
For which positions of the lines can the grid not be colored, and how few of its rays are
then needed?
\end{quote}

\begin{figure}[t]
\centering
\includegraphics[width=\textwidth]{figures/alphabet_cube_triptych.pdf}
\caption{One grid at three positions of its inner lines.  Each of the 49 rays leaves the centre
of the cube through a point where the lines cross; opposite points are the same ray.  Filled
points: the 33 rays of Peres at $x=\sqrt2$ and the 31 rays of Conway and Kochen at $x=2$.  At
$x=\sqrt3$ the whole grid of 49 rays has a coloring.  Only the two inner lines move with~$x$.
The three cubes are the points \textit{A}, \textit{B} and \textit{C} of Fig.~\ref{fig:circle}.}
\label{fig:triptych}
\end{figure}

The dial need not be real, and there need not be only one.  The same question is asked of
every grid, whatever its coordinate values, and the theorems of this paper are answers to it
for particular families of grids.

\subsection{The ladder of grids}\label{sec:ladder}

The grid of Section~\ref{sec:grid} is one of a family, and the family has an order
(Fig.~\ref{fig:ladder}).

\begin{figure}[t]
\centering
\includegraphics[width=0.84\textwidth]{figures/fig7_ladder.pdf}
\caption{Three rungs.  One face of each cube is shown, with lines at the coordinate values; the
number of rays is that of the grid.  On the first two rungs the points where the lines cross,
taken over all six faces with opposite points counted once, are all the rays of the grid.  On
the third they are 109 rays, and the grid has in general 145.}
\label{fig:ladder}
\end{figure}

On the bottom rung there is one coordinate value,~$1$.  Each face carries three lines each way
and the grid has 13 rays.  It has a coloring, so the bottom rung holds no KS set.  On the
second rung there are two coordinate values, $1$ and~$x$.  That is the grid of the two
classical sets, with 49 rays.  On the third rung there are three, $1$, $x$ and~$y$, and in
general 145 rays.  Here the drawing is no longer the whole grid.  On the cube that reaches
to~$y$, with $1<x<y$, the ray $(1,x,0)$ leaves at $(y/x,\,y,\,0)$, and $y/x$ is in general not a
coordinate value, so that point is not one where lines cross.  The crossings on the six faces
are 109 rays of the 145.  From the third rung on,
the grid is what its coordinate triples say it is, and the cube is an illustration of part of
it.

A grid that lacks some of the minus signs is a subset of one of these.  Its coloring is a
separate matter, but a KS set on it is a KS set on the larger grid, so for the question of how
few rays are needed it brings nothing new.

\paragraph{The question of the minimum.}
The best set known, of 31 rays, stands on the second rung.  Nobody has proved that 31 is the
least possible.  All that is proved in general is that a KS set in three dimensions needs at
least 24 rays~\cite{KirchwegerPeitlSzeider2023,LiBrightGanesh2024}.  The smallest set known
comes from a grid on the second rung, so the ladder poses a definite question.  Does some grid
higher up, or to the side among the complex numbers, carry a smaller set?  Or does the pattern
of the rungs show a reason why none exists?

\paragraph{What the results give.}
They do neither.  They find no smaller set, and they give no reason why there is none.
\emph{This paper does not prove that 31 is the smallest KS set in three dimensions.}  These classifications do not establish it, and the reason is simple.
Every finite set of rays lies on \emph{some} grid if enough coordinate values are allowed, so
``lies on a grid'' excludes nothing until their number is bounded.  What the rungs give is
elimination.  A KS set of fewer than 31 rays, if there is one, is not on the first rung, not
on the second for any real or complex~$x$ (Theorem~\ref{thm:twosymbol}), and not on the third
with real coordinate values (Theorem~\ref{thm:threesymbol}).  These statements are about
grids, not completed grids.  A smaller set must be sought outside them: on a grid with more
coordinate values, or with complex ones on the third rung.  The limit of this progress should
be plain.  It does not shorten a search among sets with arbitrary coordinates, because a
smaller set need not have few coordinate values.

The classification also shows something else.  On the second rung the positions at which
the grid cannot be colored are separate points.  With three coordinate values they need not
be.  For the complex grids with coordinate values $1$, $x$ and the conjugate $x^{*}$ they are
two circles, four straight lines through the origin with the origin removed, and four further
points in the plane of~$x$ (Theorem~\ref{thm:conj}).  One fixed list of 45 coordinate triples
cannot be colored at any point of the circle $|x|^2=2$ (Theorem~\ref{thm:circle}).  At a
typical point of that circle the 45 triples name 45 rays and no fewer will do.  At $x=\sqrt2$
they name only 33, the rays of Peres.  A second circle, $|x|^2=\tfrac12$, also passes through
his set, and on it the sets have 33 rays throughout (Section~\ref{sec:halfcircle}).  His set
is thus not an isolated curiosity.  It lies on two continuous families.
